% !TeX root = ../main.tex
\chapter{Related Work}
\label{ch:relatedwork}
Die Automatisierung komplexer Software-Migrationen erfordert ein tiefgreifendes Verständnis sowohl 
der proprietären Zieldomäne als auch der neuesten informationstechnischen Möglichkeiten. 
Da die wissenschaftliche Forschungslage an der direkten Schnittstelle zwischen proprietärer 
SAP-Migration und generativer Künstlicher Intelligenz bislang äußerst dünn besetzt ist, ordnet 
dieses Kapitel die vorliegende Arbeit anhand hochaktueller Publikationen aus angrenzenden 
Disziplinen in den Stand der Technik ein. Hierzu werden zunächst die spezifischen prozessualen 
Hürden bei der Transformation von SAP-Konfigurationsmodellen beleuchtet, um die Limitierungen 
aktueller Standardwerkzeuge aufzuzeigen. Daran anknüpfend wird der Einsatz von Large Language 
Models (LLMs) für die Code-Migration diskutiert, wobei insbesondere die empirisch belegten Grenzen 
monolithischer Ansätze bei strikten Architekturvorgaben herausgearbeitet werden. Abschließend wird 
der aktuelle State of the Art zu Multi-Agenten-Systemen (MAS) vorgestellt, um die methodische 
Forschungslücke für die semantische Migration von proprietärer Constraint-Logik logisch herzuleiten.

Die Migration von etablierten ERP-Systemen auf SAP S/4HANA zwingt Unternehmen zur Überführung ihrer 
historisch gewachsenen Produktmodelle von der klassischen Variantenkonfiguration (LO-VC) in die 
Advanced Variant Configuration (AVC). \textcite{matilainenChangeRequirementsProduct2024} verdeutlicht in seiner Untersuchung zu 
S/4HANA-Produktdatenstrukturen, dass dieser Systemwechsel weit über ein rein syntaktisches Update 
hinausgeht. Aufgrund der exponentiell wachsenden Variantenvielfalt bei konfigurierbaren Materialien 
wird im AVC eine strikte Modellierung durch neue Konfigurationsprofile und harte deklarative 
Constraints zwingend erforderlich.

Ein kritischer Erfolgsfaktor der Migration ist dabei die vorgelagerte Datenbereinigung. Wie 
\textcite{matilainenChangeRequirementsProduct2024} aufzeigt, müssen historische Altlasten und ungenutzte Merkmale vor der 
Systemumstellung zwingend identifiziert und entfernt werden, um die Performance und Struktur des 
neuen Systems nicht zu belasten. Obgleich SAP für diesen Prozess Standardwerkzeuge, wie 
beispielsweise Transition Workspaces und das Migration Cockpit, zur Verfügung stellt, stoßen diese 
bei der logischen Transformation an ihre Grenzen. Diese Werkzeuge arbeiten primär deterministisch 
und fokussieren sich auf die reine Massendatenübernahme. Die komplexe, semantische Restrukturierung 
von historischem Beziehungswissen in die deklarative AVC-Logik sowie die inhaltliche Bereinigung 
des Legacy-Codes verbleiben als stark fragmentierter, manueller Flaschenhals bei den Fachexperten. 
Die Literatur verdeutlicht somit, dass die konzeptionelle Vorbereitung und Logik-Transformation im 
aktuellen Industrie-Standard stark manuell geprägt ist, woraus sich ein zwingender Bedarf an 
intelligenten Automatisierungsansätzen ableitet.

Die rasanten Fortschritte im Bereich der Large Language Models (LLMs) haben die Möglichkeiten der 
automatisierten Code-Generierung und -Transformation grundlegend erweitert \parencite{karabiyikRefactorGPTChatGPTbasedMultiagent2025}. 
Dennoch belegen aktuelle Forschungsarbeiten, dass der direkte Einsatz von LLMs für die Migration 
komplexer Altsysteme – oft als „Single-Prompt-Ansatz“ oder „Monolithic Prompting“ bezeichnet – in 
realen Produktionsumgebungen kaum praktikable Ergebnisse liefert.

Die Ursache hierfür liegt primär in der Unfähigkeit monolithischer Ansätze, spezifische 
Architekturvorgaben und interne Framework-Abhängigkeiten zu berücksichtigen. \textcite{motiLegacyTranslateLLMbasedMultiAgent2026} 
demonstrieren in ihrer Untersuchung zum Framework LegacyTranslate, dass naive LLM-Übersetzungen 
bei der Migration industrieller Codebasen eine Kompilierrate von 0 \% erreichen können. Während der 
generierte Code zwar syntaktisch korrekt erscheint, scheitert er an der fehlenden Integration
notwendiger interner Schnittstellen (APIs) und domänenspezifischer Logik \parencite{motiLegacyTranslateLLMbasedMultiAgent2026}. 
Ähnliche Defizite dokumentieren Xu et al. (2025) im Kontext von MANTRA: Während spezialisierte 
Systeme hohe Erfolgsquoten erzielen, erreichen monolithische Basismodelle bei komplexen 
strukturellen Änderungen (wie dem „Move Method“-Refactoring) lediglich eine Erfolgsquote von 8,7 \%.

Darüber hinaus identifizieren \textcite{karabiyikRefactorGPTChatGPTbasedMultiagent2025} und 
\textcite{oueslatiRefAgentMultiagentLLMbased2026} konzeptionelle Schwächen 
in der Kontrollierbarkeit solcher „One-Shot“-Ansätze. Ohne eine Aufteilung in modulare 
Arbeitsschritte bleibt der Transformationsprozess eine intransparente „Blackbox“, die keine 
funktionalen Garantien bietet und bei längeren Interaktionen zu Halluzinationen neigt 
\parencite{oueslatiRefAgentMultiagentLLMbased2026,karabiyikRefactorGPTChatGPTbasedMultiagent2025}. Die Literatur macht somit deutlich, dass für eine 
erfolgreiche semantische Transformation – insbesondere unter Berücksichtigung von restriktiven 
Vorgaben wie der SAP „Clean Core“-Strategie – eine Loslösung von statischen Einzelinstruktionen hin 
zu dynamischen, mehrstufigen Prozessen zwingend erforderlich ist.

Um die Unzulänglichkeiten monolithischer Ansätze zu überwinden, rückt die aktuelle Forschung 
zunehmend Multi-Agenten-Systeme (MAS) in den Fokus, die auf dem Prinzip der verteilten Kognition 
und spezialisierter Rollenverteilung basieren. Das Kernkonzept dieser Architekturen besteht darin, 
den Transformationsprozess in logisch abgegrenzte Teilaufgaben zu zerlegen, die von kooperierenden 
Agenten autonom bearbeitet werden \parencite{oueslatiRefAgentMultiagentLLMbased2026,karabiyikRefactorGPTChatGPTbasedMultiagent2025}.

Ein wegweisendes Paradigma stellt hierbei die sequentielle Orchestrierung dar, wie sie im Framework 
RefactorGPT realisiert wird. Durch die Aufteilung in spezialisierte Rollen – von der Analyse der 
Strategie über die strukturelle Transformation bis hin zur Fehlerbehebung durch einen dedizierten 
„FixerAgent“ – wird der Migrationsprozess von einer intransparenten Blackbox in einen 
kontrollierbaren und modularen Workflow überführt \parencite{karabiyikRefactorGPTChatGPTbasedMultiagent2025}. 
Die Überlegenheit dieser Rollenspezialisierung wird durch das Framework RefAgent bestätigt: 
\textcite{oueslatiRefAgentMultiagentLLMbased2026} belegen empirisch, dass ein System aus 
spezialisierten Agenten für Planung, Generierung und Validierung die Fehlerquote und die Neigung 
zu Halluzinationen bei komplexen Interaktionen signifikant reduziert.

Ein zentrales Merkmal dieser modernen Architekturen ist die Integration autonomer Feedback-Schleifen
zur Sicherstellung der funktionalen Korrektheit. Das System MANTRA nutzt hierfür einen 
„Repair Agent“, der mittels „Verbal Reinforcement Learning“ Compiler-Fehler und Test-Logs 
analysiert, um den generierten Code iterativ zu korrigieren. Dieser Ansatz ermöglichte eine 
Steigerung der Erfolgsquote bei komplexen Transformationen von 8,7 \% auf 82,8 \% \parencite{xuMANTRAEnhancingAutomated2025}.

Speziell für die Migration in restriktiven Unternehmensumgebungen verdeutlicht das Projekt 
LegacyTranslate \textcite{motiLegacyTranslateLLMbasedMultiAgent2026} zudem die Notwendigkeit des 
sogenannten „API Grounding“. 
Da die bloße Übersetzung der Syntax in industriellen Frameworks oft zu nicht kompilierbarem Code 
führt, integriert LegacyTranslate spezialisierte Agenten, die den Zielcode aktiv an interne 
Architekturvorgaben und Bibliotheken anpassen \parencite{motiLegacyTranslateLLMbasedMultiAgent2026}. Zusammenfassend lässt der 
aktuelle Stand der Forschung erkennen, dass die Kombination aus rollenbasierter Spezialisierung, 
kontextueller Anreicherung und iterativen Validierungszyklen die einzige methodische Antwort auf 
die Komplexität industrieller Legacy-Migrationen darstellt.


Trotz der empirisch nachgewiesenen Überlegenheit von Multi-Agenten-Systemen bei der Automatisierung 
von Code-Übersetzungen offenbart die aktuelle Literatur eine signifikante Lücke hinsichtlich 
proprietärer ERP-Ökosysteme. Während etablierte Frameworks wie MANTRA \textcite{xuMANTRAEnhancingAutomated2025} 
oder LegacyTranslate \textcite{motiLegacyTranslateLLMbasedMultiAgent2026} primär auf 
objektorientierte oder klassisch prozedurale Standardsprachen (wie Java oder PL/SQL) 
fokussiert sind, fehlt bislang die Anwendung kollaborativer KI-Architekturen auf die 
domänenspezifische Constraint-Logik der SAP.

Wie die Untersuchung von \textcite{matilainenChangeRequirementsProduct2024} verdeutlicht, 
erfordert die Migration von LO-VC zu AVC nicht nur eine syntaktische Konvertierung, sondern 
eine tiefgreifende semantische Restrukturierung. Hierzu zählen insbesondere die inhaltliche 
Auflösung historischer, kognitiver Altlasten und die Überführung prozeduraler Vorbedingungen in ein 
streng deklaratives, mathematisches Lösungsnetzwerk. Bisherige wissenschaftliche MAS-Ansätze 
bieten keine architektonische Antwort auf dieses fachliche Umdenken und die strikten 
„Clean Core“-Richtlinien, die bei der SAP-Migration als harte Systemgrenzen fungieren.

Genau an dieser methodischen Schnittstelle setzt die vorliegende Arbeit an. Sie schließt die 
identifizierte Forschungslücke, indem sie die Konzepte arbeitsteiliger Agenten-Architekturen auf 
die spezifischen Herausforderungen der SAP-Variantenkonfiguration überträgt. Durch die Konzeption 
und Evaluierung eines spezialisierten Multi-Agenten-Systems zur semantischen Migration von LO-VC 
nach AVC wird ein Lösungsansatz entwickelt, der den aktuellen manuellen Flaschenhals der 
Logik-Transformation systematisch adressiert und das Systemverhalten aktiv an die moderne 
Zielarchitektur anpasst.